% 86-70(86쪽) — 좌표평면 위의 두 점 P(x축), Q(y축)와 직선 y=2x, 선분 PQ 가 직선과 만나는 점 R. 스캔 화질이 낮아 TikZ 로 다시 그림.
% 원문에 있는 것만: 직선 y=2x · 선분 PQ · 점 P·Q 의 찬 점 · 두 점이 움직이는 방향을 가리키는 화살표 · 라벨 y=2x, O, P, Q, R, x, y.
% 좌표·눈금은 원문에 없으므로 넣지 않는다(선분 OR 의 길이의 변화율이 그림에서 읽히지 않게).
\begin{tikzpicture}[x=0.70cm, y=0.70cm, line width=0.55pt]
  \draw[->, >=stealth, line width=0.7pt] (-0.70,0) -- (3.90,0) node[below=1pt, font=\small] {$x$};
  \draw[->, >=stealth, line width=0.7pt] (0,-1.15) -- (0,1.90) node[left=1pt, font=\small] {$y$};
  % 직선 y=2x
  \draw[line width=0.6pt] (-0.52,-1.04) -- (0.86,1.72);
  % 선분 PQ
  \draw[line width=0.6pt] (0,1) -- (3.05,0);
  \fill (0,1) circle (2.2pt);
  \fill (3.05,0) circle (2.2pt);
  % 움직이는 방향(원문의 화살표)
  \draw[->, >=stealth, red, line width=0.7pt] (-0.22,0.02) -- (-0.22,0.60);
  \draw[->, >=stealth, red, line width=0.7pt] (1.30,-0.26) -- (2.05,-0.26);
  \node[below right=-2pt and 0pt, font=\small] at (0,0) {$\pt{O}$};
  \node[anchor=east, inner sep=3pt, font=\small] at (0,1) {$\pt{Q}$};
  \node[below=3pt, font=\small] at (3.05,0) {$\pt{P}$};
  \node[anchor=north west, inner sep=1.5pt, font=\small] at (0.47,0.82) {$\pt{R}$};
  \node[anchor=west, font=\small] at (0.72,1.80) {$y=2x$};
\end{tikzpicture}
